\section{Introduction}
\label{sec:intro}

A whole-cell model simulates a cell from its molecular parts. The first such model
described \textit{Mycoplasma genitalium} with \pn{lit.karrsub} coupled submodels
\citep{karr2012whole}. Later models covered \textit{Escherichia coli}
\citep{macklin2020simultaneous} and the minimal cell \syn{}, first well stirred
\citep{thornburg2022fundamental} and then in space and time over a full cell cycle
\citep{thornburg2026bringing}. Each one joins
metabolism, gene expression, growth and division in one simulation, and each
reproduces measured behaviour of a real cell.

All of these models are forward simulators. Each one takes hundreds or thousands of kinetic
parameters from the literature, sets each to one value, and integrates forward in
time. The output is one trajectory. That trajectory is one sample from the
distribution of trajectories that the parameter knowledge allows. The model does not
know the width of that distribution. When a model predicts a doubling time, a reader
cannot tell whether the prediction holds for every plausible parameter set or depends
on a few uncertain values.

The missing error bar is a gap in the scientific reasoning. Every parameter in a
whole-cell model is uncertain. A rate measured in one organism, at one temperature,
in one medium, is then used in another organism under other conditions. The
uncertainty is real, and it propagates through coupled nonlinear dynamics into every
prediction. A framework that stores each parameter as a single number has no place
to put this uncertainty. It cannot ask how uncertain a prediction is, and it cannot
ask which data would reduce that uncertainty.

\subsection{Why whole-cell models do not do inference}

Bayesian inference answers both questions. It treats each parameter as a
distribution, updates the distribution with data, and propagates the result into
predictions. Systems biology argued for this approach long ago
\citep{babtie2017how}. Whole-cell models still do not use it, for the
following reasons.

\begin{enumerate}
  \item \emph{No place for uncertainty.} The frameworks store parameters as
  floating-point values in a dictionary. The source of each value is in a
  spreadsheet or a supplement, not in the model. When two sources disagree, the code
  picks one without a record.
  \item \emph{Mixed formalisms.} Metabolism is a set of ordinary differential
  equations (ODEs). Gene expression is a stochastic jump process. The two exchange
  state on a fixed clock. No single likelihood and no standard sampler covers both.
  \item \emph{Latent paths.} The stochastic block has a hidden path: the exact
  times of each transcription and translation event. The likelihood of the data
  integrates over all such paths, and that integral has no closed form.
  \item \emph{Cost.} Inference needs many simulations. For a whole-cell model the
  number of simulations, not the cost of one, sets the budget.
  \item \emph{No shared parameters.} In the published \syn{} model, no rate
  constant appears in both blocks. The blocks couple only through state. The
  standard methods for modular inference pass information through shared
  parameters \citep{plummer2015cuts,jacob2017better,carmona2020semimodular}.
  Here they have nothing to pass.
\end{enumerate}

\Cref{tab:frameworks} compares current whole-cell frameworks on these points.
Thornburg and coworkers balanced the metabolic parameters of \syn{} against each
other with a Bayesian method \citep{thornburg2022fundamental}. That step used
literature values for metabolism only. It did not infer parameters from dynamic
data through the coupled model.

\begin{table*}[t]
\caption{Capabilities of whole-cell modelling frameworks. A tick means that the
framework does this as part of its design. \placeholder{Every cell to be checked
against each framework's code and papers before submission.}}
\label{tab:frameworks}
\centering
\small
\begin{tabular}{@{}lcccc@{}}
\toprule
& Priors & Provenance & Inference & Declared \\
& & in model & on dynamics & reductions \\
\midrule
\textit{M.\ genitalium} \citep{karr2012whole} & \placeholder{?} & \placeholder{?} & \placeholder{?} & \placeholder{?} \\
\textit{E.\ coli} \citep{macklin2020simultaneous} & \placeholder{?} & \placeholder{?} & \placeholder{?} & \placeholder{?} \\
Vivarium \citep{agmon2022vivarium} & \placeholder{?} & \placeholder{?} & \placeholder{?} & \placeholder{?} \\
\syn{} \citep{thornburg2022fundamental,thornburg2026bringing} & \placeholder{?} & \placeholder{?} & \placeholder{?} & \placeholder{?} \\
\IC{} (this work) & \checkmark & \checkmark & \checkmark & \checkmark \\
\bottomrule
\end{tabular}
\end{table*}

\subsection{This work}

We present \IC{}, a framework that puts inference first. We do not claim that full
Bayesian inference over every parameter of a whole-cell model is possible today. We
claim that the framework must support inference as a basic operation, so that a
modeller can apply it where the data and the budget allow. Our contributions are:

\begin{enumerate}
  \item \textbf{A design for inference-first composition} (\cref{sec:design}). Every
  parameter is a prior with a recorded source. Modules couple through
  \pn{fw.edgekinds} kinds of typed edge. The framework checks that a composition is
  complete before it runs, and it lists every departure from the published model.
  \item \textbf{One interface with two descriptions} (\cref{sec:design}). The same
  gene-expression modules run as an exact stochastic process or as mean-field ODEs.
  In the second form the whole cell compiles to one ODE problem that automatic
  differentiation can pass through.
  \item \textbf{A reduced model of \syn{} with a real boundary} (\cref{sec:model}).
  It keeps \pn{model.rxn} metabolic reactions and \pn{model.genes} genes from the
  published model, and it keeps the two-way coupling between them.
  \item \textbf{A map of how information crosses the boundary}
  (\cref{sec:seam}), and calibrated inference across it (\cref{sec:results}). An
  exact sampler gives a calibrated reference. \placeholder{An approximate method
  agrees with it at a fraction of the cost, and scales to the full reduced model.}
\end{enumerate}

\placeholder{Last paragraph: a one-sentence statement of the main result, written
after B7 and B8.}
